% U9 WS3 -- free energy and equilibrium, dissolution, coupling. Blocks 5-6.
\documentclass{apchem-worksheet}

\apcunit{9}
\apcunittitle{Thermodynamics and Electrochemistry}
\apcdoctitle{Worksheet 3 \textbullet\ Free Energy, $K$, and Coupling}
\apcsubtitle{CED 9.5--9.7 \textbullet\ $\Delta G$ is not $\Delta G^\circ$}

\begin{document}
\wsheader[$\Delta G^\circ = -RT\ln K$ \textbullet\
$\Delta G = \Delta G^\circ + RT\ln Q$ \textbullet\
$R = \SI{8.314}{\joule\per\mole\per\kelvin}$ --- note \emph{joules}
\textbullet\ free energies of coupled steps add.]

\problem[]{\ced{9.5} Complete the correspondence:}
\begin{parts}
  \item $\Delta G^\circ$ negative $\Rightarrow$ $K$
        \answerblank[2cm]{$> 1$} $\Rightarrow$
        \answerblank[3.4cm]{products favored}
  \item $\Delta G^\circ$ zero $\Rightarrow$ $K$ \answerblank[2cm]{$= 1$}
  \item $\Delta G^\circ$ positive $\Rightarrow$ $K$
        \answerblank[2cm]{$< 1$} $\Rightarrow$
        \answerblank[3.6cm]{reactants favored}
  \item Why must $\Delta G^\circ$ be converted to joules before using
        $\Delta G^\circ = -RT\ln K$?
        \answerblank[6.4cm]{$R$ is in J/(mol$\cdot$K)}
\end{parts}

\problem[]{\ced{9.5} Reaction A has $\Delta G^\circ = -5$~kJ; reaction B
has $\Delta G^\circ = -50$~kJ. Compare their equilibrium constants
qualitatively, without a calculator.
\answerlines[3]{Both exceed 1. B's is enormously larger --- $K$ depends
exponentially on $\Delta G^\circ$, so making the free energy change ten
times more negative multiplies $K$ by a huge factor rather than by ten. A's
$K$ is only modestly above 1.}}

\problem[]{\ced{9.5} $\Delta G$ against $\Delta G^\circ$.}
\begin{parts}
  \item $\Delta G^\circ$ refers to \answerblank[4.4cm]{standard
        conditions}; $\Delta G$ refers to \answerblank[4.6cm]{the current
        mixture}.
  \item $Q < K$: sign of $\Delta G$ is \answerblank[2cm]{negative} and the
        reaction shifts \answerblank[2.4cm]{forward}.
  \item $Q = K$: $\Delta G = $ \answerblank[1.4cm]{0}
  \item $Q > K$: sign of $\Delta G$ is \answerblank[2cm]{positive}
  \item A student writes ``at equilibrium $\Delta G^\circ = 0$.'' Correct
        them. \answerlines[3]{At equilibrium $\Delta G = 0$, not
        $\Delta G^\circ$. $\Delta G^\circ$ is a fixed property of the
        reaction at that temperature and is zero only in the special case
        $K = 1$. Setting $\Delta G = 0$ in
        $\Delta G = \Delta G^\circ + RT\ln Q$ is exactly what yields
        $\Delta G^\circ = -RT\ln K$.}
\end{parts}

\problem[]{\ced{9.5} A reaction has $\Delta G^\circ = +20$~kJ. A student
concludes it can never proceed forward under any circumstances. Explain why
that is wrong.
\answerlines[3]{$\Delta G^\circ$ describes standard conditions, where every
species is at 1 M or 1 atm. If very little product is present, $Q$ is small
and $RT\ln Q$ is negative enough to make $\Delta G < 0$, so the reaction
does run forward --- until $Q$ rises to $K$.}}

\problem[]{\ced{9.5} For each, state whether the quantity is fixed at a
given temperature or depends on the current composition:}
\begin{parts}
  \item $K$: \answerblank[2cm]{fixed}
  \item $Q$: \answerblank[4.2cm]{depends on mixture}
  \item $\Delta G^\circ$: \answerblank[2cm]{fixed}
  \item $\Delta G$: \answerblank[4.2cm]{depends on mixture}
\end{parts}

\problem[]{\ced{9.6} Free energy of dissolution. Give the sign of the
contribution named:}
\begin{parts}
  \item $\Delta H$ of breaking up the lattice:
        \answerblank[2.4cm]{positive}
  \item $\Delta S$ of breaking up the lattice:
        \answerblank[2.4cm]{positive}
  \item $\Delta S$ of solvent reorganizing around the ions:
        \answerblank[2.4cm]{negative}
  \item $\Delta H$ of ion--solvent attraction:
        \answerblank[2.4cm]{negative}
\end{parts}

\problem[]{\ced{9.6} \ce{NH4NO3} dissolves endothermically --- the beaker
gets cold --- yet it is very soluble. Account for this using both terms of
$\Delta G^\circ = \Delta H^\circ - T\Delta S^\circ$.
\answerlines[4]{$\Delta H^\circ > 0$ opposes dissolution, so enthalpy alone
would predict it does not happen. But $\Delta S^\circ$ is strongly
positive: an ordered crystal becomes freely moving hydrated ions dispersed
through the solvent. At room temperature $T\Delta S^\circ$ exceeds
$\Delta H^\circ$, so $\Delta G^\circ$ is negative and the salt dissolves.
Entropy is driving this one against the enthalpy term.}}

\problem[]{\ced{9.6} The CED does not ask you to predict the overall
$\Delta G^\circ$ of dissolution from scratch. State the reason, and say
what you are expected to produce instead.
\answerlines[3]{The three contributions largely cancel --- the enthalpy
terms are large and opposite, as are the entropy terms --- so the net value
is a small difference between big numbers and is not reliably predictable
by inspection. You are expected to identify the contributions and their
signs, and reason about which dominates.}}

\problem[]{\ced{9.7} Driving an unfavorable process.}
\begin{parts}
  \item Name the two general strategies.
        \answerlines[2]{Supply external energy from outside the system, or
        couple the unfavorable reaction to a favorable one.}
  \item External energy driving electrolysis of water:
        \answerblank[2.6cm]{electrical}
  \item External energy driving photosynthesis:
        \answerblank[2cm]{light}
  \item Free energies of coupled steps \answerblank[1.6cm]{add}, provided
        the steps share a \answerblank[4.6cm]{common intermediate}.
\end{parts}

\problem[]{\ced{9.7} Phosphorylating glucose has
$\Delta G^\circ = +13.8$~kJ/mol. ATP hydrolysis has
$\Delta G^\circ = -30.5$~kJ/mol.}
\begin{parts}
  \item $\Delta G^\circ$ for the coupled process: \workspace[1.6cm]
        \keyonly{\begin{apcanswer}$+13.8 + (-30.5) =
        \mathbf{-16.7}$~kJ/mol --- favored\end{apcanswer}}
  \item Which species is the common intermediate?
        \answerblank[3cm]{\ce{P_i} (phosphate)}
  \item Explain what ``ATP is the cell's energy currency'' actually means.
        \answerlines[3]{Not that ATP stores energy in a special bond, but
        that its hydrolysis is favorable enough to pay for reactions that
        are not. Coupling the two through a shared intermediate makes the
        sum of the free energies negative, so the combined process runs.}
\end{parts}

\problem[]{\ced{9.7} Reaction X has $\Delta G^\circ = +25$~kJ; reaction Y
has $\Delta G^\circ = -18$~kJ. A student couples them and reports the pair
as favored. Evaluate.
\answerlines[3]{The sum is $+25 + (-18) = +7$~kJ, still positive, so the
coupled process is \emph{not} favored --- Y is not favorable enough to
carry X. Coupling also requires a shared common intermediate; without one
the addition is arithmetic rather than chemistry.}}

\begin{spiralstrip}{Unit 9 \textbullet\ blocks 1--4}
  \item Crossover temperature $= $ \answerblank[3cm]{$\Delta H^\circ /
        \Delta S^\circ$}
  \item $\Delta G^\circ$ for the Haber process at \SI{298}{\kelvin}:
        \answerblank[2.2cm]{$-33$ kJ}
  \item A catalyst changes: \answerblank[3cm]{$E_a$ only}
  \item $\Delta G_f^\circ$ of an element in its standard state:
        \answerblank[1.4cm]{0}
  \item $\Delta S^\circ$ for \ce{H2O(l) -> H2O(g)}:
        \answerblank[2.6cm]{$+119$ J/K}
\end{spiralstrip}

\end{document}
